% ===================================================================
% preamble.tex — Praeambel der finanzwirtschaftlichen Analyse Alamos Gold
% ===================================================================
\usepackage[T1]{fontenc}
\usepackage[utf8]{inputenc}
\usepackage{lmodern}
\usepackage[ngerman]{babel}
\usepackage{csquotes}
% patch ohne "footnote": microtype richtet das Fussnotenzeichen im Rand aus,
% kann seinen Patch aber nicht anlegen, wenn footmisc (para) den
% Fussnotenapparat ersetzt - ohne die Einschraenkung meldet jeder Lauf
% "Unable to apply patch `footnote'". Die uebrigen Patches bleiben aktiv.
\usepackage[patch={item,toc,eqnum}]{microtype}
\usepackage{xcolor}
\usepackage{amsmath}
% amssymb liefert \blacktriangle und \blacktriangledown - die Symbole, mit
% denen die Bildunterschrift von fig:kursverlauf die Markenformen benennt.
% Damit braucht die Abbildung keine Legende, die sonst in die Kurve ragt.
\usepackage{amssymb}
\usepackage{booktabs}
\usepackage{longtable}
\usepackage{array}
\usepackage{tabularx}
\usepackage{siunitx}
\usepackage{xfp}
\usepackage{graphicx}
% pgfplots bindet tikz zwar implizit ein, ein eigener Ladeaufruf macht die
% Abhaengigkeit robust gegen eine spaetere Entfernung von pgfplots.
\usepackage{tikz}
% Haelt Gleitobjekte in ihrem Abschnitt: Ohne \FloatBarrier wandert eine
% Tabelle, die den gesamten Rumpf einer Unterunterabschnitts bildet, auf die
% naechste Seite - die Ueberschrift bleibt dann ohne Inhalt zurueck.
\usepackage{placeins}
% Haelt eine Ueberschrift mit dem, was ihr folgt, zusammen. LaTeX verhindert
% von sich aus nur den Umbruch unmittelbar NACH einer Ueberschrift
% (\@afterheading setzt clubpenalty hoch); eine Ueberschrift plus eine
% einzelne Zeile am Seitenfuss bleibt erlaubt - und genau das ist zweimal
% aufgetreten. \needspace{n\baselineskip} vor der Ueberschrift schiebt sie
% auf die naechste Seite, wenn dort nicht mehr n Zeilen Platz haben.
\usepackage{needspace}
\usepackage{pgfplots}
\usepackage{pgfplotstable}
\usepackage[hidelinks]{hyperref}
\usepackage{url}

% --- Satzspiegel: Seitenraender schmaler --------------------------------
% DIV=12 ergibt bei A4 links und rechts je 74.69pt (26.25mm) Rand. Beide
% Raender werden um 10% verkuerzt (auf 67.22pt / 23.6mm); die Textbreite
% waechst dadurch von 448.13pt auf 463.07pt. Die Texthoehe (636.6pt, ein
% Vielfaches des Zeilenabstands) bleibt unveraendert, damit die
% Grundlinien und der Fussnotenapparat unangetastet bleiben.
% \areaset ist der KOMA-eigene Weg; das geometry-Paket wuerde die
% typearea-Berechnung uebersteuern.
\areaset[current]{463.07pt}{636.6pt}

% --- Schusterjungen und Hurenkinder verbieten -----------------------
% LaTeX-Vorgabe ist je 150 und damit so niedrig, dass eine einzelne Zeile
% eines Absatzes allein am Seitenfuss (Schusterjunge) oder am Seitenkopf
% (Hurenkind) stehen bleiben darf. Im deutschen Satz gelten beide als
% Fehler. 10000 verbietet sie vollstaendig; die entstehende Differenz
% faengt die Seitenhoehe auf, da KOMA hier nicht auf gleiche Satzhoehe
% zwingt. \displaywidowpenalty deckt denselben Fall vor einer abgesetzten
% Formel ab - davon gibt es in Teil 2 mehrere.
\clubpenalty=10000
\widowpenalty=10000
\displaywidowpenalty=10000

% --- Fussnotenapparat flach halten ----------------------------------
% Die Zitierweise setzt jede Quelle in eine Fussnote; untereinander gesetzt
% belegte der Apparat auf belegdichten Seiten bis zu einem Drittel des
% Satzspiegels. footmisc mit "para" setzt alle Fussnoten einer Seite in einen
% durchlaufenden Absatz ueber die volle Breite - dieselbe Ersparnis wie ein
% zweispaltiger Apparat, ohne dessen Paketkonflikt: das zweispaltige
% dblfnote (yafoot) laeuft mit pgfplots und hyperref zusammen in eine
% Endlosschleife der Ausgaberoutine (Seiten werden ohne Ende ausgegeben,
% "TeX capacity exceeded, input stack size"); reproduziert mit article und
% scrartcl, unabhaengig von der Ladereihenfolge.
\usepackage[para]{footmisc}

% --- Kopfzeile: Abschnitt links, Dokumenttitel rechts -----------------
% \sectionmark wird von jedem \section-Aufruf mit dem Titel als Argument
% aufgerufen; die Umdefinition setzt Nummer und Titel zusammen, sodass die
% Kopfzeile links "4 Beurteilung des Verlaufs des Aktienkurses" zeigt.
% \section* (Quellenverzeichnis usw.) loest \sectionmark nicht aus - die
% Kopfzeile behaelt dort den letzten nummerierten Abschnitt; deshalb wird
% sie vor dem Anhang manuell geleert (siehe hausarbeit.tex). \titlepage
% setzt bereits \thispagestyle{empty}, das Deckblatt bleibt also
% kopfzeilenfrei. Der rechte Eintrag ist statisch und benennt das Dokument.
\usepackage[automark]{scrlayer-scrpage}
\pagestyle{scrheadings}
\renewcommand{\sectionmark}[1]{\markright{\thesection\quad #1}}
\clearpairofpagestyles
% \small: Abschnittstitel plus Dokumenttitel muessen in eine Zeile passen.
\setkomafont{pagehead}{\normalfont\small}
\ihead{\rightmark}
\ohead{Analyse Alamos Gold Inc.}
\KOMAoptions{headsepline=true}
\cfoot*{\pagemark}

% --- Inhaltsverzeichnis: nur bis subsection (x.x) --------------------
% scrartcl zaehlt section=1, subsection=2, subsubsection=3. tocdepth=2
% laesst subsubsections (z. B. 5.1.2) aus dem Verzeichnis, ohne ihre
% Numerierung im Fliesstext zu aendern - secnumdepth bleibt unangetastet,
% da die Nummern der Aufgabenstellung (5.1.1-5.1.3 usw.) im Text stehen
% bleiben muessen.
\setcounter{tocdepth}{2}

\pgfplotsset{compat=1.18}
% dateplot stellt "date coordinates in=x" bereit. Ohne die Bibliothek
% versucht pgfplots, "2020-12-31" als Gleitkommazahl zu lesen und bricht ab.
\usepgfplotslibrary{dateplot}
% pgfplots formatiert seine Achsen- und Datenbeschriftungen unabhaengig von
% siunitx. Ohne diese Zeile stuenden im Diagramm Punkte, im Text Kommata.
\pgfplotsset{/pgf/number format/use comma, /pgf/number format/1000 sep={}}

% Nach unten weisendes Dreieck als eigene Markenform. Der naheliegende Weg,
% "mark=triangle*" mit "mark options={rotate=180}" zu drehen, funktioniert
% innerhalb von scatter/classes NICHT: pgfplots setzt mark options je Klasse
% neu und verwirft die Drehung stillschweigend - die Tiefstkurse wurden
% dadurch mit einem nach OBEN weisenden Dreieck gezeichnet. Eine eigene
% Markenform kann nicht ueberschrieben werden.
\pgfdeclareplotmark{dreieckab}{%
  \pgfpathmoveto{\pgfqpoint{-0.6\pgfplotmarksize}{0.5\pgfplotmarksize}}%
  \pgfpathlineto{\pgfqpoint{0.6\pgfplotmarksize}{0.5\pgfplotmarksize}}%
  \pgfpathlineto{\pgfqpoint{0pt}{-0.55\pgfplotmarksize}}%
  \pgfpathclose%
  \pgfusepathqfillstroke%
}

% Farben der drei Punktarten in fig:kursverlauf. Gruen fuer das Jahreshoch,
% Rot fuer das Jahrestief, Blau fuer den Schlusskurs. Die Farbe kommt zur
% Form hinzu und ersetzt sie nicht: Im Schwarzweissdruck bleiben die drei
% Arten ueber Kreis, Dreieck-auf und Dreieck-ab unterscheidbar.
\definecolor{kurshoch}{RGB}{0,120,60}
\definecolor{kurstief}{RGB}{175,30,35}
\definecolor{kursschluss}{RGB}{25,55,120}

% Achsenvorlage der Kursabbildung. Sie steht hier und nicht in der
% Abbildung selbst, weil scripts/check_literals.py in sections/*.tex jede
% Zahl anschlaegt und die Tickweite von 365 Tagen ein Satzparameter ist,
% keine Finanzkennzahl. Die Regel bleibt damit unangetastet.
\pgfplotsset{
  kursachse/.style={
    width=0.95\linewidth, height=6.2cm,
    date coordinates in=x,
    xticklabel style={rotate=45, anchor=north east, font=\footnotesize},
    xticklabel=\year, xtick={2021-01-01,2022-01-01,2023-01-01,2024-01-01,2025-01-01,2026-01-01},
    ylabel={Kurs in USD}, ylabel style={font=\footnotesize},
    yticklabel style={font=\footnotesize},
    ymin=0, ymax=60,
    grid=both, grid style={gray!22, very thin},
    axis lines=left, enlarge x limits=0.03},
}

% Zuordnung ueber die Spalte "art" der Datendatei (S/H/T), nicht ueber
% getrennte Datensaetze - so bleibt die Reihe eine einzige, nach Datum
% geordnete Linie.
\pgfplotsset{
  kurspunkte/.style={
    scatter, only marks, scatter src=explicit symbolic,
    scatter/classes={
      S={mark=*,         mark size=1.9pt, kursschluss},
      H={mark=triangle*, mark size=2.6pt, kurshoch},
      T={mark=dreieckab, mark size=2.6pt, kurstief}}},
}

% --- Zahlenformat: deutsches Dezimalkomma ---------------------------
\sisetup{
  output-decimal-marker = {,},
  group-separator       = {.},
  group-minimum-digits  = 4,
  % Ohne diese Zeile gruppiert siunitx auch die NACHKOMMASTELLEN: aus dem
  % Rohwert 43.4679 wurde im Satz "43,467.9" statt "43,5". Die Rohwerte in
  % data/kennzahlen.tex tragen volle Genauigkeit, deshalb trat der Fehler
  % bei jeder ungerundeten Prozentangabe auf.
  group-digits          = integer,
  detect-weight         = true,
  detect-family         = true
}

% --- Betragsmakros --------------------------------------------------
% Die Rohwerte stehen in data/kennzahlen.tex im englischen Format und in
% voller Genauigkeit. Die Formatierung erfolgt ausschliesslich hier.
%   \MioUSD       - Fliesstext, auf eine Nachkommastelle gerundet
%   \MioUSDexakt  - Rechenwege und Tabellen, volle veroeffentlichte Genauigkeit
% Der Konzernabschluss weist in Mio. USD ohne Nachkommastelle aus; die
% Rundung liegt trotzdem in der Ausgabe und nicht im gespeicherten Wert,
% damit abgeleitete Groessen (Margen, Veraenderungsraten) aus den exakten
% Werten hervorgehen und nicht aus vorgerundeten.
\newcommand{\MioUSD}[1]{\num[round-mode=places,round-precision=0]{#1}\,Mio.\,USD}
\newcommand{\MioUSDexakt}[1]{\num{#1}\,Mio.\,USD}
\newcommand{\MrdUSD}[1]{\num[round-mode=places,round-precision=2]{#1}\,Mrd.\,USD}
% Prozentangaben werden bei der AUSGABE auf eine Nachkommastelle
% gerundet; die Rohwerte bleiben exakt, damit abgeleitete Groessen
% nicht aus vorgerundeten Werten hervorgehen.
\newcommand{\Pct}[1]{\num[round-mode=places,round-precision=1]{#1}\,\%}
% Betrag einer Veraenderungsrate. Die Rohwerte in data/kennzahlen.tex tragen
% ihr Vorzeichen. In Wendungen wie "ging um ... zurueck" oder "lag um ...
% unter" traegt bereits das Verb die Richtung; das Minuszeichen waere dort
% eine doppelte Verneinung. \PctBetrag gibt denselben Wert ohne Vorzeichen aus.
\newcommand{\PctBetrag}[1]{\num[round-mode=places,round-precision=1]{\fpeval{abs(#1)}}\,\%}
\newcommand{\USDje}[1]{\num[round-mode=places,round-precision=2]{#1}\,USD}
% Faktor mit Malzeichen, z. B. "3,08x" fuer Net Debt/EBITDAC.
\newcommand{\Faktor}[1]{\num[round-mode=places,round-precision=2]{#1}\,x}
% TeX ueberliest jedes Leerzeichen NACH einem Kontrollwort. "\KurstagHochZFV
% und" wurde deshalb als "1.4.2025und" gesetzt. \Datum{} klammert das
% Datumsmakro ein und macht das Leerzeichen wieder sichtbar.
\newcommand{\Datum}[1]{#1}
% Kursangabe ohne Einheit fuer Tabellenspalten, die die Einheit im Kopf
% tragen. Zwei feste Nachkommastellen, damit 74,00 und 76,25 in derselben
% Spalte buendig stehen; ohne die feste Stellenzahl druckte siunitx "74,0".
% Unzen werden ohne Nachkommastelle und mit Tausenderpunkt gesetzt.
\newcommand{\Unzen}[1]{\num[round-mode=places,round-precision=0]{#1}\,Unzen}
\newcommand{\TsdUnzen}[1]{\num[round-mode=places,round-precision=0]{#1}\,Tsd.\,Unzen}
% Kosten und Preise je Unze. Zwei Angaben, die in diesem Geschaeft alles
% entscheiden und die deshalb ein eigenes Makro bekommen.
\newcommand{\JeUnze}[1]{\num[round-mode=places,round-precision=0]{#1}\,USD/Unze}
\newcommand{\MozReserven}[1]{\num[round-mode=places,round-precision=1]{#1}\,Mio.\,Unzen}
% Kennungen von Vorschriften und Einreichungen sind Namen und keine
% Zahlenwerte. Sie stehen hier, damit sections/*.tex frei von Ziffern bleibt
% und scripts/check_literals.py nicht abgeschwaecht werden muss.
\newcommand{\NIOffenlegung}{National Instrument~58-101}
\newcommand{\Kurswert}[1]{\num[round-mode=places,round-precision=2]{#1}}

% --- Zitierweise: deutsche Fussnotenzitation ------------------------
% Jeder Beleg ist ein Sprung ins Quellenverzeichnis: Der Kurztitel in der
% Fussnote ist mit \hyperref auf den zugehoerigen Eintrag verknuepft, ein
% Klick fuehrt dorthin. Der Verweis steht damit IN der Angabe und nicht als
% ausgeschriebene Nummer daneben.
%
% Die Farbe macht die Verknuepfung sichtbar - ohne sie waere sie zwar
% anklickbar, aber nicht erkennbar (hyperref laeuft mit "hidelinks", damit
% ringsum keine farbigen Rahmen entstehen). Fuer einen streng einfarbigen
% Ausdruck genuegt es, hier "black" einzusetzen; die Verknuepfung bleibt.
\definecolor{quellenlink}{RGB}{20,60,120}
\newcommand{\quellensprung}[2]{\hyperref[#1]{\textcolor{quellenlink}{#2}}}

% \quelleMDA{2025}{4} belegt den Lagebericht (Management's Discussion and
% Analysis) des Geschaeftsjahres 2025 auf gedruckter Seite 4. Die Seitenzahl
% ist die des gedruckten Dokuments, nicht die des HTML-Blocks der
% EDGAR-Einreichung: scripts/edgar_seiten.py liest sie aus dem Seitenfuss vor
% jedem Umbruch, und scripts/pruefe_seiten.py prueft jeden Wert dagegen.
%
% Alamos ist ein auslaendischer Emittent und reicht Form 40-F statt 10-K ein.
% Die inhaltlichen Dokumente sind deren Anhaenge: EX-99.1 Annual Information
% Form, EX-99.2 MD&A, EX-99.3 Konzernabschluss. Deshalb drei Makros und
% nicht eines.
\newcommand{\quelleMDA}[2]{\footnote{\quellensprung{quelle:mda#1}{Alamos Gold
  Inc., Management's Discussion and Analysis #1}, S.~#2.}}
\newcommand{\quelleFS}[2]{\footnote{\quellensprung{quelle:fs#1}{Alamos Gold
  Inc., Konzernabschluss #1}, S.~#2.}}
% Zwischenbericht und Circular bekommen eigene Makros statt eines freien
% Titelarguments. Grund: Der Schluessel des Quellenverzeichnisses wird aus
% dem Argument gebildet, und ein Label darf kein Makro enthalten - ein
% "\," im Titel liess den Lauf mit "Missing \endcsname inserted" abbrechen.
% Der Schluessel steht deshalb fest im Makro, die Beschriftung daneben.
\newcommand{\quelleZB}[1]{\footnote{\quellensprung{quelle:zb}{Alamos Gold
  Inc., Zwischenbericht zum 30.\,Juni 2026}, S.~#1.}}
\newcommand{\quelleZBA}[1]{\footnote{\quellensprung{quelle:zba}{Alamos Gold
  Inc., Zwischenabschluss zum 30.\,Juni 2026}, S.~#1.}}
\newcommand{\quelleCirc}[1]{\footnote{\quellensprung{quelle:circ}{Alamos Gold
  Inc., Management Information Circular 2026}, S.~#1.}}
% Kurzbeleg fuer eine Webquelle: Die Fussnote nennt Urheber und Abrufdatum und
% verweist mit dem Schluessel auf das Quellenverzeichnis, wo die vollstaendige
% Adresse steht. Das zweite Argument ist der \label-Schluessel und NICHT die
% Adresse - so steht jede URL genau einmal im Dokument, dieselbe Regel, die
% fuer Zahlenwerte gilt.
\newcommand{\quelleWeb}[3]{\footnote{\quellensprung{#2}{#1}
  (abgerufen am #3).}}

% --- Fussnotennummer wiederverwenden (manuell kuratiert) -------------
% Zitiert ein Absatz dieselbe Quelle+Seite mehrfach auf derselben
% gedruckten Seite, druckt \quelleGB/\quelleQM/\quelleWeb jedes Mal eine
% neue Fussnote mit identischem Text - auf einer Seite standen zuvor bis
% zu sieben gleichlautende Fussnoten "Aumann AG, Geschaeftsbericht 2024,
% S. 10." untereinander. \quelleMerken speichert die Nummer der gerade
% gesetzten Fussnote unter einem Schluessel; \quelleErneut verweist mit
% \footnotemark auf dieselbe Nummer, ohne den Text erneut zu drucken.
% Die betroffenen Stellen sind mit scripts/check_footnote_pages.py gegen
% die tatsaechlich gerenderte PDF-Seite verifiziert. Nach jeder Aenderung
% an Fliesstext erneut pruefen: eine verschobene Seitenumbruchgrenze kann
% eine Gruppe ungueltig machen oder eine neue entstehen lassen.
\newcommand{\quelleMerken}[1]{\expandafter\xdef\csname fn@#1\endcsname{\thefootnote}}
\newcommand{\quelleErneut}[1]{\footnotemark[\csname fn@#1\endcsname]}

% --- Drei Farben, drei Verantwortlichkeiten -------------------------
% Schwarz  Fakten, Zahlen, Tabellen, Quellen. Belegt.
% Blau     Einschaetzungen im laufenden Text, die der Verfasser prueft.
% Rot      das Votum im Verdikt-Teil: eine Empfehlung, keine Feststellung.
%
% Die Trennung ist der Kern des Dokuments. Rot ist ausdruecklich KEIN
% Beleg und keine Zusage: Es ist eine aus den vorstehenden Befunden
% abgeleitete Empfehlung unter den in Teil 5 erklaerten Annahmen. Die
% Entscheidung trifft der Verfasser, nicht das Dokument.
\definecolor{vdrot}{RGB}{155,25,30}
% \color statt \textcolor, damit mehrere Absaetze im Argument stehen
% duerfen - dieselbe Falle wie bei \bk.
\newcommand{\vd}[1]{{\color{vdrot}#1}}

% --- Blaue Zonen: vom Verfasser zu pruefende Wertungsprosa ----------
% Vor der Abgabe \newcommand{\bk}[1]{#1} setzen, um die Blaufaerbung zu entfernen.
\definecolor{bkblau}{RGB}{20,60,150}
% \textcolor ist nicht "long": ein Absatzumbruch im Argument bricht den Lauf
% mit "Paragraph ended before \@textcolor was complete" ab. Mehrere der
% blauen Zonen umfassen zwei oder drei Absaetze, deshalb \color in einer
% Gruppe statt \textcolor.
\newcommand{\bk}[1]{{\color{bkblau}#1}}

% --- Einlesen generierter Tabellenkoerper ---------------------------
% \input haengt hinter den Dateiinhalt ein \relax. Innerhalb einer
% Tabelle beginnt dieses \relax bereits die naechste Zelle, sodass das
% folgende \bottomrule mit "Misplaced \noalign" scheitert. Deshalb wird
% hier die TeX-Primitive (\@@input) verwendet, die nichts anhaengt.
% Argument ohne Endung, z. B. \tabellenkoerper{data/cf_a}.
\makeatletter
\newcommand{\tabellenkoerper}[1]{\@@input #1.tex }
\makeatother

% --- Annahmenkasten -------------------------------------------------
\newcommand{\annahme}[1]{%
  \par\medskip\noindent\fbox{\parbox{0.97\linewidth}{\small\textbf{Annahme:} #1}}\par\medskip}
